\section{Introduction}
\label{sec:intro}

JEPA-style world models predict in latent space rather than reconstruct
observations. An online encoder maps an observed context to a latent
representation, a predictor advances that latent forward, and a slowly
moving target encoder supplies the future representation to
match~\cite{grill2020bootstrap,tarvainen2017mean}. The abstraction
underlies I-JEPA~\cite{assran2023ijepa}, V-JEPA~\cite{assran2024vjepa},
and the broader JEPA programme~\cite{lecun2022path}, and it is closely
related to non-contrastive SSL methods such as BYOL~\cite{grill2020bootstrap},
SimSiam~\cite{chen2021exploring}, DINO~\cite{caron2021emerging},
VICReg~\cite{bardes2022vicreg}, and Barlow Twins~\cite{zbontar2021barlow}.
By moving away from pixel- or token-level reconstruction, JEPAs cleanly
separate representation learning from generative decoding, but they
leave open a question that becomes central under partial observability:
what geometry should the latent occupy when it must be imagined forward
without new observations~\cite{ha2018world,hafner2020dream,hafner2020dreamerv2,hafner2023dreamerv3}?

Partial observability changes the object being
predicted~\cite{kaelbling1998planning}. When the current observation
history does not determine the hidden state, the future is not a single
point in representation space. The same context supports several
compatible hidden states, and the natural latent is a belief over those
hidden variables rather than a deterministic
embedding~\cite{kaelbling1998planning,hausknecht2015drqn,igl2018deep,karkus2017qmdp}.
A vector predictor trained against an EMA target can still drive the
loss down, but the loss alone does not specify which uncertainty the
latent should preserve, which should be averaged away, or how a
counterfactual action should move the belief. We study a concrete
alternative: a JEPA whose latent is a density matrix and whose
predictor is a learned unitary on a joint system--environment
space~\cite{nielsen2010quantum,breuer2002theory,stinespring1955positive}.

The resulting model, the Unitary World Model JEPA (UWM-JEPA),
produces a joint density matrix $\rho_t$ from the online encoder and
advances it by
\begin{equation}
  \hat\rho_{t+k\mid t}
  =
  U^k \rho_t (U^\dagger)^k .
\end{equation}
The training scaffold is JEPA-native: an EMA target encoder, a
stop-gradient target, and a latent-space matching loss. Only the
geometry of the predicted quantity changes. A density matrix represents
a graded belief over latent configurations; unitary conjugation
preserves the joint-state spectrum, purity, and von Neumann entropy
exactly during blind rollout. We refer to this property as
\emph{joint-state non-forgetting}.

The qualifier \emph{joint-state} matters. The theorem applies to the
full system--environment density matrix and does not, on its own,
imply that the reduced system state remains linearly decodable after
many steps or that UWM-JEPA encodes context better than a strong
recurrent baseline. The empirical work in this paper measures what
survives the two operations that matter for downstream use: partial
trace to the system state and blind rollout through time.

\paragraph{Counterfactual targets recover an action-sensitive predictor}
Under a teacher-forced JEPA target the learned action term collapses to
$\|H_1\|/\|H_0\|\approx0.03$, because the target encoder has already
consumed the observed future and the predictor can match it without
using the action. Replacing the target by a counterfactual simulator
rollout restores the action term to $1.00\pm0.15$, action perturbations
move the predicted latent in the expected direction, and UWM-JEPA
solves a hidden-velocity indicator task on which the matched
LSTM-JEPA-CF baseline remains at majority-class accuracy.

\paragraph{The unitary predictor improves short-horizon imagination}
When the latent is rolled forward without new observations, UWM-JEPA
preserves target-nearness better than the vector predictors we tested.
In future-latent retrieval, UWM-JEPA remains far above LSTM-JEPA models
with plain and residual MLP predictors across
$k\in\{1,3,5,10\}$. In teacher-versus-blind probing, UWM-JEPA loses
less than ten points of $R^2$ at the intermediate horizons $k=1$ and
$k=3$, both shorter than the trained horizon $k=5$, while LSTM-JEPA
loses substantially more at those horizons and re-aligns near $k=5$. Beyond $k=10$, both models
have low or negative blind probe scores, so long-horizon retention
remains an open question rather than a settled claim.

\paragraph{Context representation is a control}
On the held-out linear probe used to measure context representation,
UWM-JEPA and a parameter-matched LSTM-JEPA agree within seed noise:
$R^2=0.901\pm0.032$ versus $0.894\pm0.021$ over five seeds (Welch
$p=0.70$). The density-matrix latent is not a stronger context encoder,
so the behavioral separation is not explained by encoder quality.

To our knowledge, UWM-JEPA is the first JEPA-family world model in
which the latent is treated as an explicit belief over hidden
continuations rather than a single representational point, and in
which the predictor is constrained to evolve that belief without
dissipating it during imagined rollout. The density-matrix latent,
the unitary predictor, and the counterfactual action targets are the
specific tools we use to instantiate this property, but the broader
claim is about giving a JEPA latent a place to carry uncertainty and
holding on to that uncertainty under prediction. \Cref{sec:architecture} defines the model.
\Cref{sec:nonforget} states the joint-state guarantee and its scope.
\Cref{sec:actions} gives the main action-conditioned behavioral test,
\cref{sec:imagine} tests blind rollout without new observations, and
\cref{sec:representation} establishes that these effects are not
explained by a stronger context encoder. Collapse diagnostics and the
claim-to-evidence audit are reported in the supplement.
